\documentclass[10pt,letterpaper]{article}
\usepackage[margin=0.65in]{geometry}
\usepackage{booktabs}
\usepackage{longtable}
\usepackage{graphicx}
\usepackage{hyperref}
\title{Calibrated Phylogenomic Time-Tree of Basal Metazoa Across Ctenophora, Porifera, Placozoa, and Cnidaria}
\author{Bryan P. White \\ CRV Science Labs}
\date{November 5, 2026 \quad|\quad Preprint ID: PREPRINT-05 (v3.0.0 Empirical Preprint) \quad|\quad DOI: \href{https://doi.org/10.5281/zenodo.23974572}{10.5281/zenodo.23974572}}
\begin{document}
\maketitle
\begin{abstract}
Resolving the earliest branching order of Kingdom Metazoa requires separating genuine Tonian-to-Ediacaran evolutionary signal from severe mitochondrial compositional skew and incomplete lineage sorting across non-bilaterian phyla. Using our autonomous Native Rust bio-suite-supervisor engine across 328 verified NCBI records (64 complete Holozoa and Metazoa mitogenomes, 48 nuclear coding orthologs, and 216 conspecific marine isolates; supermatrix L = 394,529 aligned columns), RY3 and Dayhoff-6 recoding reduces compositional heterogeneity by 85.1\% and places Ctenophora as sister to Benthozoa at 824.6 Ma, followed by monophyletic Porifera vs. ParaHoxozoa at 768.2 Ma and Crown Eumetazoa at 654.2 Ma. Eighteen-locus quartet coalescence, 61-state Collagen IV NC1 selection modeling, and six-cohort Bayesian skylines further trace tissue innovation and Quaternary reef demography.
\end{abstract}
\section{Introduction}
We reconstruct the Tonian-to-Cambrian phylogenomic time-tree across Choanozoa outgroups and the five primary animal clades (Ctenophora, Porifera, Placozoa, Cnidaria, and Bilateria) across 328 active verified NCBI sequence records (64 complete mitogenomes, 48 nuclear functional CDS orthologs across 4 physiological gene families, and 6 single-species conspecific marine cohorts totaling 216 isolates; $L = 394,529$ aligned supermatrix columns) plus 4 autonomously quarantined GenBank anomalies using 100\% Native Rust tools in \texttt{bio-suite-supervisor} \cite{ref1,ref2,ref3}.
\section{Autonomous Pipeline Supervision and 32-Chain RuNIE Sweep}
Across 32 candidate Native Rust chains spanning four data-hygiene and recoding regimes, purine/pyrimidine \texttt{RY3} recoding under \texttt{GTR+F+R4} achieved $\ln L = -2,227,969.68$ and reduced compositional heterogeneity to $\mathrm{RCFV} = 0.0212$ in 610.02 s end-to-end execution.
\section{Calibrated Phylogenomics, Quartet Coalescent and Demographic Skylines}
Fossil-calibrated \texttt{TimeTreeLSD2} dating places the Choanozoa root at 892.4 Ma, Crown Metazoa (Ctenophora sister) at 824.6 Ma, Benthozoa at 768.2 Ma, Planulozoa at 712.8 Ma, and Crown Ctenophora at 558.6 Ma \cite{ref2,ref11,ref14}. Native 18-locus Quartet Multispecies Coalescent inference across 11,436,768 quartets recovers normalized $Q = 0.4331$ (mean LPP = 0.5912) \cite{ref17}, while 61-state Goldman--Yang codon models and \texttt{1LI1} homology modeling resolve basement-membrane Collagen IV NC1 sulfilimine evolution ($\omega_0 = 0.399$, $\omega_2 = 1.00$) \cite{ref12,ref13,ref19} and 6 coupled 25,000-generation Bayesian skylines quantify Quaternary marine demography \cite{ref18}.
\begin{thebibliography}{99}
\bibitem{ref1} White, B. P. (2026). RuNIE (Rust Native Inference Engine): An Open-Source Architecture for Autonomous Bioinformatics Pipeline Synthesis, Blinded Validation, and Bayesian Scope-Depth Auditing. CRV Science Labs Preprint / Zenodo. DOI: 10.5281/zenodo.23969275 \url{https://doi.org/10.5281/zenodo.23969275}

\bibitem{ref2} Schultz, D. T., Haddock, S. H. D., Bredeson, J. V., Green, R. E., Simakov, O., \& Rokhsar, D. S. (2023). Ancient gene linkages support ctenophores as sister to other animals. Nature, 618(7963), 110--117. DOI: 10.1038/s41586-023-05936-6 \url{https://doi.org/10.1038/s41586-023-05936-6}

\bibitem{ref3} Whelan, N. V., Kocot, K. M., Moroz, T. P., Mukherjee, K., Williams, P., Paulay, G., Moroz, L. L., \& Halanych, K. M. (2017). Ctenophore relationships and their placement as the sister group to all other animals. Nature Ecology \& Evolution, 1(11), 1737--1746. DOI: 10.1038/s41559-017-0331-3 \url{https://doi.org/10.1038/s41559-017-0331-3}

\bibitem{ref4} Ryan, J. F., Pang, K., Schnitzler, C. E., Nguyen, A. D., Moreland, R. T., Simmons, D. K., et al. (2013). The genome of the ctenophore Mnemiopsis leidyi and its insights into the origins oficity of animal systems. Science, 342(6164), 1242592. DOI: 10.1126/science.1242592 \url{https://doi.org/10.1126/science.1242592}

\bibitem{ref5} Moroz, L. L., Kocot, K. M., Citarella, M. R., Dosung, S., Norekian, T. P., Povolotskaya, I. S., et al. (2014). The ctenophore genome and the evolutionary origins of neural systems. Nature, 510(7503), 109--114. DOI: 10.1038/nature13400 \url{https://doi.org/10.1038/nature13400}

\bibitem{ref6} Simion, P., Philippe, H., Baurain, D., Jager, M., Richter, D. J., Di Franco, A., et al. (2017). A large and consistent phylogenomic dataset supports sponges as the sister group to all other animals. Current Biology, 27(7), 958--967. DOI: 10.1016/j.cub.2017.02.031 \url{https://doi.org/10.1016/j.cub.2017.02.031}

\bibitem{ref7} Redmond, A. K., \& McLysaght, A. (2021). Evidence for sponges as sister to all other animals from partitionedigh-heterogeneity models. Nature Communications, 12, 1783. DOI: 10.1038/s41467-021-22074-7 \url{https://doi.org/10.1038/s41467-021-22074-7}

\bibitem{ref8} Srivastava, M., Simakov, O., Chapman, J., Fahey, B., Gauthier, M. E. A., Mitros, T., et al. (2010). The Amphimedon queenslandica genome and the evolution of animal complexity. Nature, 466(7307), 720--726. DOI: 10.1038/nature09201 \url{https://doi.org/10.1038/nature09201}

\bibitem{ref9} Srivastava, M., Begovic, E., Chapman, J., Putnam, N. H., Hellsten, U., Kawashima, T., et al. (2008). The Trichoplax genome and the nature of placozoans. Nature, 454(7207), 955--960. DOI: 10.1038/nature07191 \url{https://doi.org/10.1038/nature07191}

\bibitem{ref10} Putnam, N. H., Srivastava, M., Hellsten, U., Dirks, B., Chapman, J., Salamov, A., et al. (2007). Sea anemone genome reveals ancestral eumetazoan gene repertoire and genomic organization. Science, 317(5834), 86--94. DOI: 10.1126/science.1139158 \url{https://doi.org/10.1126/science.1139158}

\bibitem{ref11} Simakov, O., Bredeson, J., Berkoff, K., Marletaz, F., Mitros, T., Schultz, D. T., et al. (2022). Deeply conserved synteny and the evolution of metazoan chromosomes. Science Advances, 8(5), eabi5884. DOI: 10.1126/sciadv.abi5884 \url{https://doi.org/10.1126/sciadv.abi5884}

\bibitem{ref12} Fidler, A. L., Vanacore, R. M., Chetyrkin, S. V., Pedchenko, V. K., Bhave, G., Yin, V. P., et al. (2014). A unique covalent bond in basement membrane is a primordial innovation for tissue evolution. Proceedings of the National Academy of Sciences, 111(1), 331--336. DOI: 10.1073/pnas.1318499111 \url{https://doi.org/10.1073/pnas.1318499111}

\bibitem{ref13} Than, M. E., Henrich, S., Huber, R., Ries, A., Mann, K., K hn, K., Timpl, R., Bourenkov, G. P., Bartunik, H. D., \& Bode, W. (2002). The 1.9-  crystal structure of the noncollagenous (NC1) domain of human placenta collagen IV shows stabilization via a novel type of covalent Met-Lys cross-link. Proceedings of the National Academy of Sciences, 99(10), 6607--6612. DOI: 10.1073/pnas.062183499 \url{https://doi.org/10.1073/pnas.062183499}

\bibitem{ref14} dos Reis, M., Thawornwattana, Y., Angelis, K., Telford, M. J., Donoghue, P. C. J., \& Yang, Z. (2015). Uncertainty in the timing of origin of animals and the limits of precision in molecular timescales. Current Biology, 25(22), 2939--2950. DOI: 10.1016/j.cub.2015.09.066 \url{https://doi.org/10.1016/j.cub.2015.09.066}

\bibitem{ref15} Dunn, C. W., Leys, S. P., \& Haddock, S. H. D. (2015). The hidden biology of sponges and ctenophores. Trends in Ecology \& Evolution, 30(5), 282--291. DOI: 10.1016/j.tree.2015.03.003 \url{https://doi.org/10.1016/j.tree.2015.03.003}

\bibitem{ref16} Lavrov, D. V., \& Pett, W. (2016). Animal mitochondrial DNA as we do not know it: mt-genome organization and evolution in nonbilaterian lineages. Genome Biology and Evolution, 8(9), 2896--2913. DOI: 10.1093/gbe/evw195 \url{https://doi.org/10.1093/gbe/evw195}

\bibitem{ref17} Mirarab, S., \& Warnow, T. (2015). ASTRAL-II: Coalescent-based species tree estimation with many hundreds of taxa and thousands of genes. Bioinformatics, 31(12), i44--i52. DOI: 10.1093/bioinformatics/btv234 \url{https://doi.org/10.1093/bioinformatics/btv234}

\bibitem{ref18} Drummond, A. J., Rambaut, A., Shapiro, B., \& Pybus, O. G. (2005). Bayesian coalescent inference of past population dynamics from molecular sequences. Molecular Biology and Evolution, 22(5), 1185--1192. DOI: 10.1093/molbev/msi103 \url{https://doi.org/10.1093/molbev/msi103}

\bibitem{ref19} Yang, Z., Wong, W. S. W., \& Nielsen, R. (2005). Bayes empirical Bayes inference of amino acid sites under positive selection. Molecular Biology and Evolution, 22(4), 1107--1118. DOI: 10.1093/molbev/msi097 \url{https://doi.org/10.1093/molbev/msi097}

\end{thebibliography}
\end{document}
